\section{Simultaneous concentration in lattice quotients}\label{sec:cube}

Retaining prescribed layers forces a new bit column into an intermediate
lattice. We now bound the number of bit columns in every coset of every
such lattice, on one typical event for the starting matrix.

Let \(B\) be a nonsingular binary matrix of size \(m\in\mathcal I_N\), and
let
\[
 B\Z^m\subseteq K\subseteq\Z^m.
\]
For \(D\subset[m]\), put
\[
 Q_D=\{z\in\bits^m:z_j=0\text{ for }j\notin D\},\qquad
 M_K(D)=\max_{a\in\Z^m}|Q_D\cap(a+K)|,
 \qquad M_K=M_K([m]).
\]
In particular, \(M_K(D')\le M_K(D)\) whenever \(D'\subset D\).
The lattice \(K\) has full rank, so \(\det K=[\Z^m:K]\). Define
\[
 q=q_K=\frac{\log[\Z^m:K]}{N\log N},
 \qquad\text{so that}\qquad [\Z^m:K]=N^{qN}.
\]
The quotient \(\Z^m/K\) is a quotient of \(G_B\), so
Equation~\eqref{eq:group-order} gives \(0\le q\le1/2\).

For \(0\le d\le1\), write \(f=1-d\) and define
\begin{equation}\label{eq:cube-profile}
 \begin{split}
 c(d/2)&=1-fz(d/f),\\
 z(u)&=
 \begin{cases}
 (1+\sqrt{2u-u^2})^{-1},&0\le u\le1,\\
 1/2,&u\ge1.
 \end{cases}
 \end{split}
\end{equation}
At \(f=0\), the product \(fz(d/f)\) is defined to be zero. Thus
\(c:[0,1/2]\to[0,1]\), with \(c(0)=0\) and \(c(1/2)=1\).

\begin{proposition}[Cube concentration]\label{prop:cube-concentration}
For every fixed \(A>0\), typically for nonsingular binary matrices \(B\)
of size \(m\in\mathcal I_N\), one has simultaneously for every lattice
\(B\Z^m\subseteq K\subseteq\Z^m\)
\begin{equation}\label{eq:cube-concentration}
 M_K\le 2^{N(1-c(q_K)+\ell^{-A})}.
\end{equation}
The typical event may depend on \(A\), but it does not depend on \(K\).
Thus the choice of \(K\) may depend arbitrarily on \(B\) and on any
other data.
\end{proposition}

The function \(c\) is nondecreasing and concave. To see this, put
\(r(u)=\sqrt{2u-u^2}\) for \(0<u<1\). Then
\[
 r'(u)=\frac{1-u}{r(u)},\qquad r''(u)=-\frac1{r(u)^3}.
\]
Consequently \(z=(1+r)^{-1}\) is nonincreasing and convex on \([0,1]\);
more explicitly,
\[
 z''(u)=\frac{1}{r(u)^3(1+r(u))^2}
       +\frac{2(1-u)^2}{r(u)^2(1+r(u))^3}>0.
\]
Its left derivative at \(1\) is zero, so its constant extension remains
convex and nonincreasing. The perspective \(fz(d/f)\) is convex for
\(f>0\): this follows by applying convexity of \(z\) with weights
proportional to the two values of \(f\). Its continuous extension to
the endpoint \(f=0\) is also convex. Restriction to \(f=1-d\) proves
concavity of \(c\). As \(d\) increases, both \(f\) and \(z(d/f)\) are
nonnegative and nonincreasing, which proves monotonicity.

We will use the uniform bound
\begin{equation}\label{eq:cube-holder}
 |c(q)-c(q')|\le C_{\mathrm{pr}}|q-q'|^{1/2}
 \qquad(0\le q,q'\le1/2)
\end{equation}
for an absolute constant \(C_{\mathrm{pr}}\). Indeed, for \(u,v\in[0,1]\),
\[
 |z(u)-z(v)|\le |r(u)-r(v)|
 \le \sqrt{2|u-v|}.
\]
On \(0\le d\le1/2\), the map \(d\mapsto d/(1-d)\) is \(4\)-Lipschitz,
so the product in Equation~\eqref{eq:cube-profile} is
Hölder of exponent \(1/2\) there. On \(1/2\le d\le1\) it is
\((1-d)/2\). Combining the two bounds at \(d=1/2\), and then using
\(d=2q\), gives Equation~\eqref{eq:cube-holder}.

We first count bases of low-height subspaces, and then use that count
to price the witnesses that will arise from a large cube intersection.

\subsection{Low-height bases and witness pricing}

\begin{lemma}[Exact-basis count]\label{lem:exact-basis-count}
Fix \(C,C'>0\). Uniformly for \(m\in\mathcal I_N\) and \(0\le j\le m\),
the number of ordered tuples
\[
 (y,y_1,\ldots,y_j)\in(\bits^m)^{j+1}
\]
for which \(y_1-y,\ldots,y_j-y\) are independent and their span \(W\)
satisfies \(H(W)\le\exp(N\ell^C)\) is at most
\[
 2^{j^2+o(N^2\ell^{-C'})}.
\]
\end{lemma}

\begin{proof}
The case \(j=0\) is immediate. For \(j\ge1\), encode the center \(y\)
and a choice of \(j\) pivot coordinates on which the differences remain
independent. These choices cost \(O(N\log N)\) bits. At the pivots,
each difference entry has two possible values once \(y\) is known, so
all pivot entries cost \(j^2\) bits.

Add the other coordinates in a fixed order. Suppose the current
projected span is \(V\subset\R^h\), and put \(L_V=V\cap\Z^h\). The next
projected span is the graph of a rational functional
\(x\mapsto\langle v,x\rangle\) on \(V\), represented by \(v\in V\). Put
\[
 a=\min\{b\in\Z_{>0}:bv\in L_V^*\}.
\]
The image of \(L_V\) under this functional modulo \(\Z\) has order
\(a\). Hence the projection of the integer lattice in the graph has
index \(a\) in \(L_V\). The volume factor of the graph map on \(V\) is
\(\sqrt{1+\|v\|^2}\), and therefore
\begin{equation}\label{eq:cube-height-ratio}
 \frac{H(\operatorname{graph}(v))}{H(V)}
       =a\sqrt{1+\|v\|^2}.
\end{equation}
All these ratios are at least one. The initial projected space is
\(\R^j\), of height one, and the final one is \(W\). Thus every
intermediate height is at most \(\exp(N\ell^C)\), and a ratio exceeding
\(N^{1/10}\) occurs at most \(O(N\ell^C/\log N)\) times. Mark those
coordinates and encode their \(j\) difference entries directly. Their
total cost is \(O(N^2\ell^C/\log N)\) bits.

At any other coordinate,
\[
 a\le N^{1/10},\qquad \|av\|\le N^{1/10}.
\]
We count possible \(av\) in \(\Lambda=L_V^*\).
Lemma~\ref{lem:normal-height}, applied to \(V\subset\R^h\), gives
\(\det M\ge H(V)^{-1}\) for every sublattice \(M\subset\Lambda\).

Apply Lemma~\ref{lem:det-minimizer} to \(\Lambda\) with
\(T=N^{-1/10}\), choosing a minimizing sublattice as a function of the
known \(V\). If its rank is \(r_0\), comparison with the zero lattice
and the preceding determinant bound give
\[
 H(V)^{-1}\le T^{r_0},\qquad
 r_0\le \frac{10N\ell^C}{\log N}.
\]
Every sublattice of the quotient has determinant at least \(T\) to
its rank. By Lemma~\ref{lem:gaussian-count}, the number of quotient
images of vectors of norm at most \(N^{1/10}\) is
\[
 2^{O(N^{2/5}\log^2N)}.
\]
Here the quotient has rank at most \(N\), and the bound also covers
rank zero.

Encode \(a\), costing \(O(\log N)\) bits, and the quotient image of
\(av\). In a fixed fiber, \(v\) varies in an affine space of dimension
at most \(r_0\). Evaluation on the already known \(j\) independent
projected differences is injective on \(V\), and the evaluations are
the \(j\) entries of the new difference column. Each entry has two
possible values, determined by the corresponding coordinate of \(y\).
An affine space of dimension \(r_0\) meets such a cube in at most
\(2^{r_0}\) points, by projection onto \(r_0\) suitable coordinates.
This bounds the fiber cost by \(r_0\) bits.

Summing over at most \(N\) coordinates, including the pivot, center,
and marking costs, bounds the logarithm of the number of tuples by
\[
 j^2+O\!\left(
   \frac{N^2\ell^C}{\log N}+N^{7/5}\log^2N+N\log N
 \right).
\]
For fixed \(C,C'\), the error is \(o(N^2\ell^{-C'})\), uniformly in
\(j\) and \(m\), as required.
\end{proof}

\begin{proposition}[Witness pricing]\label{prop:witness-pricing}
Fix \(D_0>0\), and put \(\eta=\ell^{-D_0}\) and \(\rho=\eta^3\).
Typically for nonsingular binary matrices \(B\) of size
\(m\in\mathcal I_N\), the following assertion holds simultaneously
for every \(B\Z^m\subseteq K\subseteq\Z^m\).

Let \(S_1,\ldots,S_h\subset\bits^m\) be nonempty subsets, each
\(\rho\)-robust on its affine span, where \(h\le3/\eta\). Suppose
\(S_i-S_i\subset K\), and put \(W_i=\Span(S_i-S_i)\). Suppose also
that \(W=\sum_i W_i\) is a direct sum. Write
\[
 \dim W_i=t_iN,\qquad t=\sum_i t_i,\qquad |S_i|=2^{\alpha_iN}.
\]
If \(D\subset[m]\) has \(|D|=xN\), \(x\le t\), and
\(M_K(D)\le2^{\eta N}\), then
\begin{equation}\label{eq:witness-pricing}
 xt-\frac{x^2}{2}
 \le \sum_i t_i(t_i-\alpha_i)+\eta t+\eta^3.
\end{equation}
\end{proposition}

For a fixed nonsingular matrix \(B\), a \emph{witness} consists of data
\(K,(S_i)_{i=1}^h,D\) satisfying the stated assumptions on the lattice,
sets, their difference spans, and \(D\). It is \emph{violating} if
Equation~\eqref{eq:witness-pricing} fails.
The two sides of that equation have a counting interpretation. In a
symmetric block on \(tN\) coordinates, the number
of independent entries with at least one endpoint in a specified set
of \(xN\) coordinates is
\((xt-x^2/2)N^2+O(N)\). Restricting each row on those coordinates to
at most \(2^{\eta N}\) choices saves this many bits, less
\(\eta tN^2\). For the \(i\)-th set \(S_i\), describing an ordered basis
costs \(t_i^2N^2\) bits to leading order, while robustness supplies
\(\alpha_i t_iN^2\) bits of basis multiplicity. The term
\(t_i(t_i-\alpha_i)\) is the remaining cost. A violation would make
the row saving exceed the total witness cost by \(\eta^3N^2\).

\begin{proof}
The robust-height bound in Lemma~\ref{lem:robust-height}, with the
fixed parameter \(\rho=\ell^{-3D_0}\), gives
\[
 H(W_i)\le\exp(O(N\ell)).
\]
The integer lattice sum of \(W_i\cap\Z^m\) is contained in
\(W\cap\Z^m\). The volume inequality for this sum and saturation give
\[
 H(W)\le\prod_i H(W_i)
       \le\exp(O(N\ell/\eta)).
\]
Choose a fixed \(D_1>D_0+1\). For sufficiently large \(N\), this is at
most \(\exp(N\ell^{D_1})\). We impose
Proposition~\ref{prop:typical-ranks} at this height and at a fixed
precision strong enough that its rank error is at most \(\eta^5N\).
We also impose Proposition~\ref{prop:prime-corank}. Their intersection
is typical. All counting below takes place within this intersection.

For a fixed list of dimensions, Lemma~\ref{lem:exact-basis-count}
counts centers and ordered bases of differences for the pieces with
total logarithmic cost
\begin{equation}\label{eq:witness-basis-upper}
 N^2\sum_i t_i^2+o(N^2\eta^4).
\end{equation}
Indeed, use that Lemma with a fixed \(C'\) larger than \(5D_0\) and
sum its errors over at most \(3/\eta\) pieces. The height exponent
there can be chosen fixed and large enough for all the \(W_i\).

For each individual witness, fix a center in each \(S_i\). At every
stage of choosing an ordered basis of differences, the affine span of
the center and the points already chosen is a proper affine
subspace of \(\aff S_i\). Robustness bounds the fraction of \(S_i\)
in that subspace by \(2^{-\rho}\). Consequently the number of choices
of an ordered basis for this piece is at least
\[
 \bigl((1-2^{-\rho})\,2^{\alpha_iN}\bigr)^{t_iN}.
\]
The choices for different pieces form distinct basis lists. As
\(t\le m/N\le1\), their total logarithm is at least
\begin{equation}\label{eq:witness-basis-lower}
 N^2\sum_i t_i\alpha_i
       -O\bigl(N\log(1/\rho)\bigr)
 =N^2\sum_i t_i\alpha_i-o(N^2\eta^4).
\end{equation}

Fix basis-list data, so in particular \(W\) is known. For a witnessed
matrix choose a coordinate set \(Y_0\) of size \(tN\) containing \(D\).
The second rank assertion of Proposition~\ref{prop:typical-ranks}
shows that the columns on \([m]\setminus Y_0\), projected off \(W\),
have rank at least \(m-tN-\eta^5N\). Choose an independent collection
there, and use nonsingularity of \(B\) to complete it to a basis of
the quotient \(\R^m/W\) using columns of \(B\). Let \(P\) be the
resulting set of \(m-tN\) indices and \(Y=[m]\setminus P\). Then
\[
 |D\cap P|\le\eta^5N.
\]
Encode \(D\) and \(P\), at cost \(O(N)\) bits, and reveal the columns
on \(P\). If \(T_m=m(m+1)/2\) is the total number of independent bits
of \(B\), these columns cost
\[
 T_m-\binom{tN+1}{2}
\]
bits, using symmetry.

The revealed columns and the basis differences are \(m\) independent
integer vectors: the differences form a basis of \(W\), and the
revealed columns project to a basis off \(W\). Let \(L\) be their
integer span. It is known from the encoded data, satisfies
\(L\subset K\), and has full rank. Each generating vector has norm
at most \(\sqrt N\), so Hadamard's inequality gives
\[
 |\Z^m/L|\le N^{m/2}\le N^{N/2}.
\]
We next enumerate the possible \(K\) containing this known \(L\).
Let \(r_N=\lceil N^{3/5}\rceil\). For every prime \(p\),
\[
 \dim_{\F_p}\bigl((\Z^m/K)/p(\Z^m/K)\bigr)
 \le\dim_{\F_p}(G_B/pG_B)
 =\corank_{\F_p}B\le r_N
\]
by Proposition~\ref{prop:prime-corank}. The minimum number of
generators of a finite abelian group is the maximum of these
dimensions over its primes, as follows from its invariant factors.
Thus \(\Z^m/K\), and also its character dual, can be generated by at most
\(r_N\) elements.

Put \(H=\Z^m/L\). The annihilator of \(K/L\) identifies
\((\Z^m/K)^\vee\) with a subgroup of
\(H^\vee=\Hom(H,\mathbb Q/\Z)\); by finite duality this subgroup determines
\(K\). Every subgroup with at most
\(r_N\) generators is generated by some ordered \(r_N\)-tuple in
\(H^\vee\), padding with zero elements if necessary. The number of
possible \(K\) is therefore at most \(|H|^{r_N}\). Its logarithm is
\begin{equation}\label{eq:witness-lattice-cost}
 O(N^{8/5}\log N)=o(N^2\eta^4).
\end{equation}
This candidate list is formed after the basis lists and revealed
columns fix \(L\), before the remaining matrix block is encoded.
It may include lattices incompatible with every completion. The
prime-corank condition is used only to ensure that each actual
witnessed \(K\) appears in the list.

Write \(D'=D\cap Y\), \(y=|Y|=tN\), and \(d'=|D'|\). We have
\(d'\ge(x-\eta^5)N\) and \(M_K(D')\le M_K(D)\le2^{\eta N}\). In the
remaining \(Y\)-block, first encode the symmetric bits on
\(Y\setminus D'\), at cost \(\binom{y-d'+1}{2}\). For a row indexed
by \(Y\setminus D'\), all entries outside \(D'\) are now known.
The full row belongs to \(K\), since \(B\) is symmetric and each
column belongs to \(B\Z^m\). Its entries on \(D'\) therefore lie in a
fixed coset of \(K\) and have at most \(M_K(D')\) possibilities.
This encodes the cross block at cost at most \(\eta N(y-d')\).
For the final \(D'\)-block, the same bound applies to each of its
\(d'\) rows. Ignoring symmetry within that block only enlarges the
count, so its cost is at most \(\eta Nd'\). Relative to an unrestricted
symmetric \(Y\)-block the saving is at least
\[
 \begin{split}
 \binom{y+1}{2}-\binom{y-d'+1}{2}-\eta Ny
 &=yd'-\frac{(d')^2}{2}+\frac{d'}2-\eta Ny\\
 &\ge N^2\left(xt-\frac{x^2}{2}-\eta t
                  -O(\eta^5+1/N)\right).
 \end{split}
\]

It remains to compare the upper count with the multiplicity of
witnesses. Group violating witnesses by \(h\), the integers \(t_iN\)
and \(xN\), and the integers
\[
 a_i=\lfloor\log_2|S_i|\rfloor,\qquad \bar\alpha_i=a_i/N.
\]
There are at most \(2^{O(\eta^{-1}\log N)}\) groups. Fix one group,
and let \(\mathcal E\) be the set of matrices in the imposed typical
intersection that have a violating witness in this group. For each
\(B\in\mathcal E\), select one such witness and count the distinct
pairs consisting of \(B\) and the basis lists supplied by that
witness. Equation~\eqref{eq:witness-basis-lower} gives a lower bound
of
\[
 |\mathcal E|\,
 2^{N^2\sum_i t_i\bar\alpha_i-o(N^2\eta^4)}
\]
on the number of pairs. On the other hand,
Equations~\eqref{eq:witness-basis-upper} and
\eqref{eq:witness-lattice-cost}, together with the matrix saving,
give the upper bound
\[
 2^{T_m+
 N^2\left(\sum_i t_i^2-xt+x^2/2+\eta t\right)
 +o(N^2\eta^4)}.
\]
The enumeration may count a pair more than once if it has several
compatible choices of \(D,P,K\); this preserves the upper bound.
In particular, the lower bound selects just one witness per matrix,
while the upper bound includes the union over all \(K\).

Since \(0\le\alpha_i-\bar\alpha_i<1/N\) and \(t\le1\), violation of
Equation~\eqref{eq:witness-pricing} implies
\[
 \sum_i t_i(t_i-\bar\alpha_i)-xt+x^2/2+\eta t
 <-\eta^3+\frac1N.
\]
All binary matrices have probability \(2^{-T_m}\). Dividing the pair
bounds therefore gives
\[
 \Pr(\mathbf B_m\in\mathcal E)
 \le 2^{-N^2\eta^3+o(N^2\eta^4)}.
\]
The group count and all omitted costs are absorbed in this error.
Since \(N\eta^3\to\infty\), the total probability of a violation is
\(2^{-\omega(N)}\), including the exceptional probabilities of the
imposed rank and corank properties. This proves the assertion, including its
simultaneity over the lattices, lists, and coordinate sets.
\end{proof}

\subsection{From witness pricing to concentration}

\begin{proof}[Proof of Proposition~\ref{prop:cube-concentration}]
Fix \(A>0\). Choose a fixed \(D_0>2A\), and put
\(\eta=\ell^{-D_0}\), \(\rho=\eta^3\). Work on the
typical event in Proposition~\ref{prop:witness-pricing}, and fix an
arbitrary \(K\) with \(B\Z^m\subseteq K\subseteq\Z^m\).

Apply Lemma~\ref{lem:det-minimizer} to \(K\) with
\(T=(\log N)^4\), and let \(K_0\) be a minimizing saturated sublattice.
Put \(F=\Span K_0\) and
\[
 d=\frac{m-\dim F}{N},\qquad
 f=1-d=\frac{\dim F+N-m}{N}.
\]
The projected quotient \(K/K_0\) has rank \(dN\) and all its
sublattices have determinant at least \(T\) to their rank. Moreover
\(\det K_0\le T^{\dim F}\), by comparison with the zero lattice.
The projections of the columns of \(B\) span \(F^\perp\), and each
has norm at most \(\sqrt N\). Choosing \(dN\) independent projections
and using that their integer span is a sublattice of \(K/K_0\) gives
\(\det(K/K_0)\le N^{dN/2}\). Hence
\begin{equation}\label{eq:cube-order-dimension}
 q=\frac{\log\det K}{N\log N}
 \le \frac d2+\frac{\dim F}{N}\frac{\log T}{\log N}
 \le \frac d2+\frac{4\ell}{\log N}.
\end{equation}

Consider any nonempty subset \(S\) of a coset of \(K\) in the cube, and base
its differences at one point of \(S\). Their projections off \(F\)
belong to \(K/K_0\) and have norm at most \(\sqrt N\).
Lemma~\ref{lem:gaussian-count} gives at most
\[
 2^{O(N/\log^6N)}
\]
distinct projections. Restrict to a largest projection fiber and
then apply Lemma~\ref{lem:robust-extraction} with parameter
\(\rho=\eta^3\). For sufficiently large \(N\), these two operations
lose at most \(2\eta^3N\) bits of logarithmic size. Thus, if
\(|S|=2^{aN}\), they give a \(\rho\)-robust subset \(S'\) with
\begin{equation}\label{eq:cube-fiber-extraction}
 |S'|\ge2^{(a-2\eta^3)N},
 \qquad \Span(S'-S')\subset F.
\end{equation}
If \(S\) is supported on a coordinate set, the same is true of its
differences after these operations.

Choose a set \(P_F\) of \(\dim F\) pivot coordinates for \(F\), so
coordinate projection onto \(P_F\) is injective on \(F\), and put
\(J_0=[m]\setminus P_F\). Then \(|J_0|=dN\) and a vector supported on
\(J_0\) belongs to \(F\) only if it is zero. Distinct bit vectors
supported on \(J_0\) therefore have distinct projections off \(F\):
equal projections would give a difference in \(F\). Applying the
preceding projection count to a coset intersection on \(J_0\) gives
\begin{equation}\label{eq:cube-pivot-complement}
 M_K(J_0)\le2^{O(N/\log^6N)}\le2^{\eta N}.
\end{equation}

Write \(M_K=2^{\alpha N}\). We have \(0\le\alpha\le1\). If
\(\alpha<10\eta\), the conclusion follows for sufficiently large
\(N\), since \(1-c(q)\ge0\) and \(10\eta\le\ell^{-A}\). We may
therefore assume \(\alpha\ge10\eta\).

Apply Equation~\eqref{eq:cube-fiber-extraction} to a full-coordinate
coset attaining \(M_K\). Denote the resulting robust subset by \(S_*\),
put \(W_*=\Span(S_*-S_*)\), and write
\[
 |S_*|=2^{\alpha_*N},\qquad v=\frac{\dim W_*}{N}.
\]
An affine space of dimension \(\dim W_*\) contains at most
\(2^{\dim W_*}\) cube points, by coordinate projection. Consequently
\begin{equation}\label{eq:cube-full-witness}
 \alpha_*\ge\alpha-2\eta^3,\qquad
 \alpha-2\eta^3\le v\le\frac{\dim F}{N}\le f.
\end{equation}
The last inequality uses \(m\le N\).

We construct one coordinate set \(D\) with small \(M_K(D)\) and use it
in two applications of Proposition~\ref{prop:witness-pricing}. A list of
robust sets will force \(D\) to be large; applying the estimate to the
full-coordinate set \(S_*\) then restricts how large \(\alpha\) can be
for that same \(D\).

Initially all coordinates of \(P_F\) are available. At step \(i\),
choose, among the available coordinates, a set \(A_i\) of
\emph{minimum cardinality} such that
\[
 M_K(J_0\cup A_i)>2^{\eta N}.
\]
For a coset attaining \(M_K(J_0\cup A_i)\), apply
Equation~\eqref{eq:cube-fiber-extraction}. It gives a robust subset
\(S_i\) whose difference span satisfies
\[
 W_i\subset F\cap E_{J_0\cup A_i},\qquad
 \alpha_i=\frac{\log_2|S_i|}{N}\ge\eta-2\eta^3,\qquad
 t_i=\frac{\dim W_i}{N}\ge\alpha_i.
\]
Because projection of \(F\) onto \(P_F\) is injective and \(W_i\)
vanishes on \(P_F\setminus A_i\), we have \(t_iN\le|A_i|\). Choose
\(t_iN\) pivot coordinates \(H_i\subset A_i\) on which projection of
\(W_i\) is injective, and declare those coordinates unavailable.
Stop when \(t=\sum_i t_i\ge\alpha-2\eta\).

The required set \(A_i\) exists before the stopping condition is met.
Indeed, the coordinates already declared unavailable number \(tN\).
In a full coset attaining \(M_K\), fix their bit values to a most
frequent pattern. At least \(2^{(\alpha-t)N}>2^{2\eta N}\) points
remain. Subtracting the fixed pattern on those coordinates translates
the coset and leaves bit vectors supported on \(J_0\) and the available
coordinates. Hence that coordinate set has \(M_K>2^{\eta N}\).
Equation~\eqref{eq:cube-pivot-complement} also shows that each \(A_i\)
is nonempty.

The spaces \(W_i\) form a direct sum. In a relation \(\sum_i w_i=0\)
with \(w_i\in W_i\), the later terms vanish on the pivots \(H_1\);
injectivity on \(W_1\) gives \(w_1=0\), and the same argument
successively gives every \(w_i=0\). Thus \(t\le\dim F/N\le1\).
Each increment is at least \(\eta-2\eta^3\), so there are fewer than
\(3/\eta\) steps for sufficiently large \(N\). Together with the
existence just proved, this ensures that the stopping condition is
reached. The family of available subsets only shrinks from one step
to the next. Since each \(A_i\) has minimum cardinality, the numbers
\(|A_i|\) do not decrease. If the last one has size \(sN\), then
\begin{equation}\label{eq:cube-coordinate-list}
 t\ge\alpha-2\eta,\qquad t_i\le |A_i|/N\le s
 \quad\text{for every }i.
\end{equation}

Remove one coordinate from the last \(A_i\) and include all of \(J_0\).
The resulting set \(D\) satisfies
\begin{equation}\label{eq:cube-small-coordinate-set}
 |D|=(d+s)N-1,\qquad M_K(D)\le2^{\eta N},
\end{equation}
by the minimum cardinality of that \(A_i\). For \(x\ge0\), \(t>0\),
define
\[
 e(x,t)=
 \begin{cases}
 x-x^2/(2t),&0\le x\le t,\\
 t/2,&x\ge t.
 \end{cases}
\]
This function is nondecreasing in both variables. Its rate of change
in \(x\) is between zero and one, and its rate of change in \(t\) is
between zero and \(1/2\). These assertions follow from
\[
 \partial_x e=1-x/t,\quad \partial_t e=x^2/(2t^2)\quad(x<t),
 \qquad
 \partial_x e=0,\quad \partial_t e=1/2\quad(x>t),
\]
and continuity at \(x=t\).

If necessary, truncate \(D\) to a subset of size \(tN\).
Monotonicity of \(M_K(D)\) preserves its bound, so
Proposition~\ref{prop:witness-pricing} applied to the constructed list
gives
\[
 \begin{split}
 e(d+s,t)
 &\le \frac1t\sum_i t_i(t_i-\alpha_i)
       +\eta+\frac{\eta^3}{t}+\frac1N\\
 &\le s+O(\eta).
 \end{split}
\]
The \(1/N\) accounts for the coordinate removed in
Equation~\eqref{eq:cube-small-coordinate-set}. The last inequality
uses \(\alpha_i\ge0\), \(t_i\le s\), and
\(t\ge\alpha-2\eta\ge8\eta\). Replacing \(t\) by \(\alpha\) keeps
this bound: if \(t\ge\alpha\), use monotonicity in the second
variable; if \(t<\alpha\), use its \(1/2\)-Lipschitz bound and
\(\alpha-t\le2\eta\). Thus \(e(d+s,\alpha)\le s+O(\eta)\).

Put \(h=d+s\). If \(h\le\alpha\), the last inequality gives
\[
 d\le \frac{h^2}{2\alpha}+O(\eta).
\]
Since \(\alpha\le1\), it follows that
\(h\ge\sqrt{2d\alpha}-O(\sqrt\eta)\). If \(h\ge\alpha\), it instead
gives \(s\ge\alpha/2-O(\eta)\), and then
\[
 h\ge d+\alpha/2-O(\eta)
   \ge\sqrt{2d\alpha}-O(\eta).
\]
We have proved
\begin{equation}\label{eq:cube-coordinate-lower}
 d+s\ge\sqrt{2d\alpha}-O(\sqrt\eta).
\end{equation}

Use the same \(D\) again, now with the singleton list consisting of
\(S_*\). Truncate \(D\) to size \(vN\) if necessary. From
Proposition~\ref{prop:witness-pricing} and
Equation~\eqref{eq:cube-full-witness},
\[
 e(d+s,v)
 \le v-\alpha_*+\eta+\frac{\eta^3}{v}+\frac1N
 \le v-\alpha+O(\eta).
\]
Here \(v\ge\alpha-2\eta^3\ge9\eta\) for large \(N\). Since \(v\le f\)
and the rate of change of \(e\) in its second argument is at most
\(1/2\), this yields
\begin{equation}\label{eq:cube-coordinate-upper}
 e(d+s,f)
 \le e(d+s,v)+\frac{f-v}{2}
 \le f-\alpha+O(\eta).
\end{equation}

We finish by solving the two coordinate inequalities. In the present
case \(f\ge v>0\). If \(d\ge f\), then \(d+s\ge f\), so
Equation~\eqref{eq:cube-coordinate-upper} gives
\[
 \alpha\le f/2+O(\eta)=fz(d/f)+O(\eta).
\]
Suppose instead that \(d\le f\). Monotonicity and the unit Lipschitz
bound of \(e\) in its first argument, together with
Equations~\eqref{eq:cube-coordinate-lower} and
\eqref{eq:cube-coordinate-upper}, give
\begin{equation}\label{eq:cube-profile-inequality}
 e(\sqrt{2d\alpha},f)\le f-\alpha+O(\sqrt\eta).
\end{equation}
Put \(u=d/f\in[0,1]\). Directly from
Equation~\eqref{eq:cube-profile},
\[
 2uz(u)\le1,\qquad
 1-z(u)=\sqrt{2uz(u)}-uz(u).
\]
For completeness, if \(r=\sqrt{2u-u^2}\), then
\((1-u)^2z(u)^2-2z(u)+1=0\), which gives the second identity after
taking the nonnegative square root. The first is equivalent to
\(2u\le1+r\); for \(u\le1/2\) it is immediate, and for
\(u\ge1/2\) its square follows from
\(r^2-(2u-1)^2=(1-u)(5u-1)\ge0\). These identities show that
\(\alpha_0=fz(u)\) satisfies
\[
 e(\sqrt{2d\alpha_0},f)=f-\alpha_0.
\]
The function \(a\mapsto a+e(\sqrt{2da},f)\) is nondecreasing and
increases by at least the increase in \(a\). Therefore
Equation~\eqref{eq:cube-profile-inequality} implies
\(\alpha\le\alpha_0+O(\sqrt\eta)\). Both cases give
\[
 \alpha\le fz(d/f)+O(\sqrt\eta)
          =1-c(d/2)+O(\sqrt\eta).
\]

By Equation~\eqref{eq:cube-order-dimension}, monotonicity of \(c\),
and Equation~\eqref{eq:cube-holder},
\[
 c(q)\le c(d/2)+O\!\left(\sqrt{\frac{\ell}{\log N}}\right).
\]
It follows that
\[
 \alpha\le 1-c(q)+O(\sqrt\eta)
                    +O\!\left(\sqrt{\frac{\ell}{\log N}}\right).
\]
Because \(D_0>2A\) is fixed, both error terms are
\(o(\ell^{-A})\). Increasing \(N\) makes their sum at most
\(\ell^{-A}\), proving Equation~\eqref{eq:cube-concentration}.
The argument was deterministic for every \(K\) on the single typical
event of Proposition~\ref{prop:witness-pricing}, so the conclusion
is simultaneous in \(K\).
\end{proof}
